% F0 -- önceden yazılmış karar kurallarıyla protokol akışı (şematik; veri içermez). Şekil içi metin ve başlık Türkçe.
\begin{figure*}[t]
\centering
\resizebox{\textwidth}{!}{%
\begin{tikzpicture}[
  box/.style={rectangle, rounded corners=3pt, draw=black!70, fill=black!5, line width=0.8pt,
    align=center, text width=3.9cm, minimum height=1.1cm, inner sep=4pt, font=\footnotesize},
  gate/.style={box, draw=iyiRenk, fill=iyiRenk!10, line width=1.3pt},
  stop/.style={box, draw=kotuRenk, fill=kotuRenk!8, line width=1.3pt},
  lock/.style={box, draw=black, fill=black!15, line width=2pt},
  arr/.style={-{Latex[length=2.6mm,width=2mm]}, line width=1.1pt, black!80},
  lab/.style={font=\fontsize{6}{7}\selectfont, fill=white, inner sep=1pt, align=center}
]
% ---- satır 1: taban model geçerliliği ve kart
\node[box,  minimum height=1.5cm] (A) at (0,0)    {\textbf{Replikasyon kapısı}\\3 taban model tohumu, VOC07 test\\literatüre karşı $\pm$ tolerans};
\node[gate, minimum height=1.5cm] (B) at (5.0,0)  {\textbf{Kapı durumu?}\\YEŞİL / SARI / KIRMIZI};
\node[box,  minimum height=1.9cm] (C) at (10.0,0) {\textbf{Hipotez kartı}\\koşumdan önce dondurulur:\\tek değişken, mekanizma,\\beklenen $\Delta$, alarm düzeyi,\\kapatma kuralı};
\node[stop, minimum height=0.9cm] (S) at (5.0,-2.0) {SARI: protokol açığını kaydet\\KIRMIZI: dur ve teşhis et};
% ---- satır 2: aşamalı koşum (sağdan sola)
\node[box] (D) at (10.0,-3.5) {\textbf{Aşama 1: duman testi}\\verinin \%5--10'u, 2--3 epoch};
\node[box] (E) at (5.0,-3.5)  {\textbf{Aşama 2: kısa ön deneme}\\tüm veri, 5 epoch};
\node[box] (F) at (0,-3.5)    {\textbf{Aşama 3: tam koşum}\\100 epoch, tohum 0};
% ---- satır 3: kararlar
\node[gate, minimum height=1.5cm] (G) at (0,-7.0)    {\textbf{Ön-eleme}\\$\Delta_{s0}\ge+0.30$ puan?};
\node[box,  minimum height=1.5cm] (H) at (5.0,-7.0)  {\textbf{Tohum 1 ve 2}\\taban model tohumlarıyla eşli;\\turda en fazla bir kart};
\node[gate, minimum height=1.9cm] (I) at (10.0,-7.0) {\textbf{SÜRDÜR kuralı?}\\ortalama $\Delta\ge+0.40$,\\tek yönlü eşli $p<0.10$,\\$\ge2$ tohumda $\Delta>0$,\\alarm yok};
% ---- satır 4: sonuçlar
\node[stop, minimum height=1.5cm] (K) at (5.0,-10.5)  {\textbf{Kapat ya da incele}\\kart kapatılır, negatif\\sonuç olarak raporlanır};
\node[lock, minimum height=1.9cm] (J) at (10.0,-10.5) {\textbf{Kilitli test, tek açılış}\\dondurulmuş plan; taban model,\\aday kartlar, eğitimsiz satır;\\her biri 3 tohum};
% ---- oklar
\draw[arr] (A) -- (B);
\draw[arr] (B) -- (C) node[lab, pos=0.5, above=2pt] {YEŞİL};
\draw[arr] (B) -- (S);
\draw[arr] (C) -- (D);
\draw[arr] (D) -- (E) node[lab, pos=0.5, above=2pt] {kapı};
\draw[arr] (E) -- (F) node[lab, pos=0.5, above=2pt] {kapı};
\draw[arr] (F) -- (G);
\draw[arr] (G) -- (H) node[lab, pos=0.5, above=2pt] {evet};
\draw[arr] (H) -- (I);
\draw[arr] (G.south) |- (K.west);
\node[lab] at (0,-9.2) {hayır};
\draw[arr] (I.south west) -- (K.north east) node[lab, pos=0.5] {hayır};
\draw[arr] (I) -- (J) node[lab, pos=0.5] {evet};
\draw[arr] (K) -- (J) node[lab, pos=0.5, above=2pt, align=center] {dur\\kuralı};
\end{tikzpicture}}%
\caption{Önceden yazılmış karar kurallarıyla analiz protokolünün akışı. Mavi çerçeveli kutular karar kapısı (koşumdan önce yazılmış), kırmızılar durdurma veya kart kapanışı, kalın çerçeve kilitli testtir.}\label{fig:f0-protokol}
\end{figure*}
